% Appendix A: A Bernoulli Error-Theory Primer
% DEFINES: app:bernoulli; sec:bernoulli-model, sec:bernoulli-axioms,
%   sec:bernoulli-cipher-link
% REFERENCES (backward): def:correctness, thm:composition, prop:space-duality
% Referenced by: Ch 4 (eta = Bernoulli FNR), Ch 11 (FPR/FNR accounting).

\chapter{A Bernoulli Error-Theory Primer}
\label{app:bernoulli}

The book uses, at several points, an error model it does not develop: the
correctness parameter $\eta$ is identified with a false-negative rate, errors
are asserted to compose multiplicatively, and the space bound
$-\log_2\varepsilon + H(Y)$ is attributed to an information-theoretic
argument. All three come from the Bernoulli model of approximate sets and
maps, a separate family of papers. This appendix gives the minimum needed to
make the book self-contained on those points and points to the papers for
the rest.

\section{The Approximate-Set Model}
\label{sec:bernoulli-model}

A \emph{Bernoulli set} is an approximate membership test: a structure that,
queried with an element, answers ``in'' or ``out'' correctly most of the
time. Two error rates characterize it. The \emph{false-positive rate}
$\alpha$ is the probability that a non-member is reported as a member; the
\emph{false-negative rate} $\beta$ is the probability that a member is
reported as a non-member. Together they are a $2 \times 2$ confusion matrix,
a noisy channel from the latent membership bit to the reported one:
\[
  \begin{array}{c|cc}
     & \text{reported in} & \text{reported out} \\\hline
    \text{truly in}  & 1 - \beta & \beta \\
    \text{truly out} & \alpha    & 1 - \alpha
  \end{array}
\]
A Bloom filter is the familiar instance, with $\beta = 0$ (no false
negatives) and a tunable $\alpha$; the hash-based cipher map of
\Cref{ch:hash-constructions} is another, and the two error rates are exactly
the cipher map's parameters seen from the error side.

\section{The Two Axioms}
\label{sec:bernoulli-axioms}

What makes the model tractable, and what the book's composition results rest
on, is that the errors of different elements are independent. The Bernoulli
model assumes two things: that each element's error is statistically
independent of every other's (element-wise independence), and that the error
rates do not depend on which other elements were queried (conditional
independence of block error rates). Under these axioms the joint behavior of
$m$ elements does not require an exponential joint confusion matrix; it
factors.

\begin{quote}
\textbf{Kronecker factorization.} Under the two axioms, the joint confusion
matrix for $m$ elements is the tensor (Kronecker) product of the $m$
per-element channels.
\end{quote}

The consequence is the one the book uses repeatedly: a quantity that would
otherwise need $\Theta(2^m)$ joint parameters needs only the $2$ per element.
This is why composition is predictable rather than intractable, and it is the
formal content behind the Bernoulli Axiom named among the book's core
principles.

\section{The Connection to Cipher Maps}
\label{sec:bernoulli-cipher-link}

Three of the book's results are Bernoulli results in cipher-map clothing.

First, \emph{correctness is a false-negative rate}. For the batch hash
construction of \Cref{ch:hash-constructions}, the correctness parameter
$\eta$ of \Cref{def:correctness} is exactly the Bernoulli false-negative
rate $\beta$: the fraction of in-domain elements whose hash does not land in
the correct codeword region. The error budget the seed search trades against
is the FNR of the underlying approximate set.

Second, \emph{composition multiplies}. The composition theorem
\Cref{thm:composition}, $\eta_{\mathrm{total}} \le 1 - \prod_i(1 - \eta_i)$,
is the Bernoulli composition closure theorem: under element-wise
independence, the probability that a chain of $m$ approximate tests all
succeed is $\prod_i(1 - \eta_i)$, and the chain errs with the complementary
probability. The inequality, rather than equality, is the book's own
refinement (a stage error can be masked downstream); the multiplicative
survival core is the Bernoulli model's.

Third, \emph{space meets an entropy bound}. The duality
\Cref{prop:space-duality}, $-\log_2\varepsilon + H(Y)$ bits per element, is
derived in the Bernoulli entropy analysis as an information-theoretic lower
bound and achieved by the Bernoulli hash-function construction: the
$-\log_2\varepsilon$ term is the cost of separating signal from a noise floor
of density $\varepsilon$, and the $H(Y)$ term is the source-coding cost of
the value distribution.

\section{Further Reading}
\label{sec:bernoulli-further}

The full theory, false-positive and false-negative accounting, the confusion
matrix algebra, the composition closure theorem, the entropy and space
bounds, and the extension to approximate maps and algebraic types, is
developed in the Bernoulli paper family \cite{towell2026bernoulli}: the set
model, the composition theorem for set operations, the approximate-map model,
the space-optimal hash construction, and the entropy and space complexity of
random approximate sets. This appendix has taken from it only the three
results the book leans on; the rest stands on its own.
